% Copy this folder, replace metadata/body/bibliography, then compile with latexmk -xelatex.
\documentclass[UTF8,a4paper,fontset=fandol,zihao=5,openany]{ctexrep}
\usepackage[margin=25.4mm]{geometry}
\usepackage{graphicx,booktabs,amsmath}
\usepackage[sort&compress]{gbt7714}
\usepackage{caption}
\usepackage[hidelinks]{hyperref}
\IfFontExistsTF{Times New Roman}{\setmainfont{Times New Roman}}{\setmainfont{TeX Gyre Termes}}
% Fandol is bundled with TeX Live: deterministic CJK choice on Windows and macOS.
\ctexset{
  chapter={format=\centering\heiti\bfseries\fontsize{16}{20}\selectfont,
    beforeskip=12pt,afterskip=6pt},
  section={format=\heiti\bfseries\fontsize{14}{18}\selectfont},
  subsection={format=\heiti\bfseries\fontsize{12}{16}\selectfont}}
\setlength{\parindent}{2em}
\setlength{\parskip}{6pt}
\captionsetup{position=bottom,font=small,labelsep=space}
\renewcommand{\bibfont}{\fontsize{10}{12}\selectfont}
\setcounter{tocdepth}{2}

% Replace these sample values before report delivery; do not invent personal details.
\newcommand{\ReportCourse}{课程名称}
\newcommand{\ReportTitle}{报告题目}
\newcommand{\ReportCollege}{学院系别}
\newcommand{\ReportMajor}{专业名称}
\newcommand{\ReportStudent}{学生姓名}
\newcommand{\ReportStudentID}{学生学号}
\newcommand{\ReportTeacher}{任课教师}
\newcommand{\ReportDate}{提交日期}

\begin{document}
\fontsize{10.5}{13.125}\selectfont
\pagestyle{empty}
\begin{titlepage}
  \centering
  \vspace*{22mm}
  {\heiti\bfseries\fontsize{18}{24}\selectfont 《\ReportCourse》\par}
  \vspace{12mm}
  {\heiti\bfseries\fontsize{22}{28}\selectfont 课程报告\par}
  \vspace{22mm}
  {\fontsize{13}{22}\selectfont
  \begin{tabular}{rl}
    实验题目 & \underline{\makebox[90mm]{\ReportTitle}}\\[6mm]
    学院系别 & \underline{\makebox[90mm]{\ReportCollege}}\\[6mm]
    专业名称 & \underline{\makebox[90mm]{\ReportMajor}}\\[6mm]
    学生姓名 & \underline{\makebox[90mm]{\ReportStudent}}\\[6mm]
    学生学号 & \underline{\makebox[90mm]{\ReportStudentID}}\\[6mm]
    任课教师 & \underline{\makebox[90mm]{\ReportTeacher}}
  \end{tabular}\par}
  \vfill
  {\ReportDate\par}
  \vspace*{10mm}
\end{titlepage}
\pagenumbering{gobble}
\tableofcontents
\thispagestyle{empty}
\clearpage
\pagenumbering{arabic}
\pagestyle{plain}

\chapter{研究目的}
% Replace all instructional sample content with the actual report.
本文为课程报告模板示例。正式报告应说明研究问题、任务要求和评价标准，所有结果须有可核验的来源。
\section{任务与范围}
围绕课程要求组织正文，区分已观察结果、解释和待验证假设。在线技术文档可作为工具使用说明的来源\cite{typst-docs}。

\chapter{方法与结果}
\section{方法说明}
公式示例：
\begin{equation}\label{eq:mean}
\bar{x}=\frac{1}{n}\sum_{i=1}^{n}x_i
\end{equation}
式\eqref{eq:mean}表示算术平均值；实际分析应说明数据来源和适用条件。
\begin{figure}[htbp]
  \centering
  \fbox{\parbox{.72\linewidth}{\centering 研究问题 $\longrightarrow$ 方法设计 $\longrightarrow$ 证据分析}}
  \caption{研究过程示意}\label{fig:process}
\end{figure}
如图\ref{fig:process}所示，结论应由证据支持。
\begin{table}[htbp]
  \centering
  \begin{tabular}{ll}\toprule
  项目 & 记录要求\\\midrule
  输入 & 保留数据来源\\
  输出 & 保存实际结果\\\bottomrule
  \end{tabular}
  \caption{证据记录要求}\label{tab:evidence}
\end{table}

\chapter{讨论与结论}
结合实际结果说明限制，避免将模板示例作为实验结论。表\ref{tab:evidence}仅展示记录结构。

\clearpage
\phantomsection
\addcontentsline{toc}{chapter}{参考文献}
\bibliographystyle{gbt7714-numerical}
\bibliography{references}
\end{document}
